\section{Future Work}
\label{app:future}

GPU-CFR fixes the inner loop of tabular CFR and leaves the algorithm that calls it
untouched, so the natural next steps are the solvers that call that loop many times.
Real-time subgame re-solving \cite{burch2014solving,li2026real} runs a fresh solve
under a per-decision budget on every turn; a captured graph turns each of those
solves into one launch, and the budget then buys iterations instead of dispatch.
Abstraction pipelines \cite{li2026effective,li2026abstraction} score candidate
abstractions by solving each one, and computing tournament continuations in place of
the ICM heuristic \cite{li2026icm} solves one subtree per stack configuration; both
are loops of short, structurally identical solves that the compile-once design serves
directly. On the evaluation side, the exact best-response and value-decomposition
passes that variance-reduced estimators need \cite{li2026correlated,li2026av,li2026agents}
share the compiled arrays of \cref{sec:compiled}. Finally, solver-guided language
agents \cite{li2026pokerskill,wang2026solver} query equilibrium strategies during
reasoning, and a single-GPU solver that answers those queries in milliseconds removes
the latency that today limits such agents to precomputed charts. Extending the compiler
to more than two players and to sampled (Monte Carlo) CFR variants, whose traversal is
no longer static, is the main open design question.
